% ═══════════════════════════════════════════════════════════════════════════════
% CAPÍTULO 6 — RETROPROPAGACIÓN: EL MÉTODO ADJUNTO SOBRE EL GRAFO COMPUTACIONAL
% ═══════════════════════════════════════════════════════════════════════════════
% Cambios principales respecto a la versión anterior:
%   1. Apertura reconectada con Cap. 5: "tenemos el optimizador, falta el
%      gradiente eficiente".
%   2. Ejemplo del oscilador: BP manual sobre el MLP (1,32,1) del Cap. 4,
%      valores numéricos concretos, verificación con diferencias finitas.
%   3. Ejemplos del método adjunto: reemplazados por ejemplos universales
%      para cualquier físico (mecánica analítica, óptica, termodinámica,
%      problemas inversos en espectroscopía). Se mantiene Tarantola/sismología
%      por su importancia histórica.
%   4. Algoritmo forward+backward preservado como Algoritmo central.
%   5. Notación Γ del Cap. 5 consistente.
%   6. Estrella Polar independiente al final.
%   7. Transición al Cap. 7 explícita.
% ═══════════════════════════════════════════════════════════════════════════════

\chapter{Retropropagación: el método adjunto sobre el grafo computacional}
\label{ch:cap6}

% ─────────────────────────────────────────────────────────────────────────────
\section{El problema del gradiente eficiente}
\label{sec:prob_grad}
% ─────────────────────────────────────────────────────────────────────────────

\dificultad{2}{2}

El Capítulo~\ref{ch:cap5} construyó una jerarquía de algoritmos de
optimización —GD, SGD, momento, Adam— que comparten un ingrediente: el
gradiente $\nabla_{\bm{\theta}}\mathcal{L} \in \mathbb{R}^p$. Ese
gradiente se usó en todas las reglas de actualización sin explicar cómo
calcularlo eficientemente. Esta deuda se salda ahora.

El problema es de escala. El MLP $(1, 32, 1)$ del oscilador tiene
$p = 97$ parámetros. Un modelo de lenguaje moderno tiene $p \sim 10^{10}$.
La aproximación más directa al gradiente —diferencias finitas— requiere
$p + 1$ evaluaciones del forward pass: una por cada parámetro perturbado.
Para $p = 10^{10}$ y un forward pass de $1\,\text{ms}$, el costo sería
$\sim 10^7$ segundos ($\sim 4$ meses) por iteración. Imposible.

\begin{proposicion}{Costo de las diferencias finitas}
\label{prop:costo_df}
Para un modelo con $p$ parámetros, la estimación del gradiente por
diferencias finitas centradas:
\[
  \frac{\partial\mathcal{L}}{\partial\theta_j}
  \approx \frac{\mathcal{L}(\bm{\theta}+h\bm{e}_j)
  - \mathcal{L}(\bm{\theta}-h\bm{e}_j)}{2h}
\]
requiere $2p$ evaluaciones del forward pass. El error de truncación es
$O(h^2)$ y el error de redondeo es
$O(\varepsilon_{\text{máquina}}/h)$, con balance óptimo en
$h \sim \varepsilon_{\text{máquina}}^{1/3} \approx 10^{-5}$.
\end{proposicion}

El resultado central de este capítulo es que
$\nabla_{\bm{\theta}}\mathcal{L}$ puede calcularse con un forward pass
más un backward pass, cada uno de costo comparable al forward solo. El
costo total es $\sim 3\times$ un forward pass, independiente de $p$. Este
resultado no es un truco específico de redes neuronales: es la aplicación
del \textbf{método adjunto} al grafo computacional, una técnica que el
físico computacional conoce bajo distintos nombres —multiplicadores de
Lagrange, principio de Pontryagin, operador adjunto— y que aquí aparece
en su forma más pura.

% ─────────────────────────────────────────────────────────────────────────────
\section{Diferenciación automática: dos modos}
\label{sec:autodiff}
% ─────────────────────────────────────────────────────────────────────────────

\dificultad{2}{2}

La diferenciación automática (AD) calcula derivadas exactas (salvo
redondeo) de forma algorítmica, sin derivar expresiones simbólicas ni
usar diferencias finitas. Existen dos modos.

\textbf{Modo forward (tangente):} se propaga una perturbación
$\delta\bm{\theta}$ a través del grafo, calculando
$\delta\mathcal{L} = \nabla\mathcal{L}^\top\delta\bm{\theta}$ (una
derivada direccional) en un solo paso. Para obtener el gradiente
completo se necesitan $p$ pasadas con
$\delta\bm{\theta} = \bm{e}_1, \ldots, \bm{e}_p$: costo total $O(p)$
forward passes.

\textbf{Modo reverso (adjunto):} se propaga el gradiente de la salida
$\bar{\mathcal{L}} = 1$ hacia atrás, calculando
$\nabla_{\bm{\theta}}\mathcal{L}$ completo en un solo paso. Costo:
$O(1)$ backward passes, independiente de $p$.

Para una función $f: \mathbb{R}^p \to \mathbb{R}$ (pérdida escalar con
muchos parámetros), el modo reverso es exponencialmente más eficiente que
el forward. Backpropagation es el modo reverso aplicado al MLP.

El modo forward es preferible cuando la función tiene pocas entradas y
muchas salidas ($f: \mathbb{R}^p \to \mathbb{R}^m$ con $p \ll m$): el
jacobiano $\bm{J} \in \mathbb{R}^{m \times p}$ se calcula en $p$
pasadas forward (cada una da una columna de $\bm{J}$) vs.\ $m$ pasadas
reverso (cada una da una fila). La Estrella Polar de este capítulo
contiene un ejemplo donde el modo forward es la elección correcta.

% ─────────────────────────────────────────────────────────────────────────────
\section{La regla de la cadena en grafos computacionales}
\label{sec:cadena_grafo}
% ─────────────────────────────────────────────────────────────────────────────

\dificultad{2}{2}

\subsection{Composición y regla de la cadena matricial}
\label{subsec:cadena}

Para $\mathcal{L} = g(f(\bm{\theta}))$ con
$f: \mathbb{R}^p \to \mathbb{R}^m$ y $g: \mathbb{R}^m \to \mathbb{R}$:
\begin{equation}
  \nabla_{\bm{\theta}}\mathcal{L}
  = \bm{J}_f(\bm{\theta})^\top\,\nabla_{\bm{z}}g(f(\bm{\theta})),
  \label{eq:cadena}
\end{equation}
donde $\bm{J}_f \in \mathbb{R}^{m \times p}$ es el jacobiano de $f$.
Para una composición de $L$ funciones
$\mathcal{L} = g^{(L)} \circ \cdots \circ g^{(1)}(\bm{\theta})$:
\begin{equation}
  \nabla_{\bm{\theta}}\mathcal{L}
  = \bm{J}_{g^{(1)}}^\top \bm{J}_{g^{(2)}}^\top \cdots
  \bm{J}_{g^{(L)}}^\top \cdot 1.
  \label{eq:cadena_comp}
\end{equation}
El modo reverso evalúa este producto de derecha a izquierda: primero
$\bm{J}_{g^{(L)}}^\top \cdot 1$ (vector en $\mathbb{R}^{n_{L-1}}$),
luego $\bm{J}_{g^{(L-1)}}^\top$ aplicado al resultado, y así
sucesivamente. Cada paso es un producto jacobiano-transpuesto-por-vector,
nunca la formación explícita del jacobiano.

\subsection{Gradiente adjunto y regla local}
\label{subsec:adjunto_local}

Para un nodo del grafo con entrada $\bm{x}$ y salida
$\bm{y} = f(\bm{x})$, se define:
\begin{align}
  \bar{\bm{y}} &= \partial\mathcal{L}/\partial\bm{y}
  \quad\text{(gradiente adjunto de la salida, recibido del nodo siguiente)},
  \label{eq:adjunto_salida}\\
  \bar{\bm{x}} &= \partial\mathcal{L}/\partial\bm{x}
  \quad\text{(gradiente adjunto de la entrada, enviado al nodo anterior)}.
  \label{eq:adjunto_entrada}
\end{align}
La regla de backpropagation local es:
\begin{equation}
  \bar{\bm{x}} = \bm{J}_f(\bm{x})^\top\bar{\bm{y}}.
  \label{eq:bp_local}
\end{equation}
El jacobiano transpuesto es el \textbf{operador adjunto} de la
transformación forward: el gradiente fluye hacia atrás por la traspuesta.
Para una capa afín $\bm{y} = \bm{W}\bm{x}$, el adjunto es
$\bar{\bm{x}} = \bm{W}^\top\bar{\bm{y}}$.

% ─────────────────────────────────────────────────────────────────────────────
\section{Las cuatro ecuaciones de backpropagation}
\label{sec:cuatro_ec}
% ─────────────────────────────────────────────────────────────────────────────

\dificultad{3}{3}

\subsection{Gradientes de operaciones elementales}
\label{subsec:bp_elem}

\begin{proposicion}{Backpropagation de la transformación afín}
\label{prop:bp_afin}
Para $\bm{z} = \bm{W}\bm{a} + \bm{b}$ con
$\bar{\bm{z}} = \partial\mathcal{L}/\partial\bm{z}$:
\begin{equation}
  \bar{\bm{a}} = \bm{W}^\top\bar{\bm{z}}, \qquad
  \partial\mathcal{L}/\partial\bm{W} = \bar{\bm{z}}\bm{a}^\top, \qquad
  \partial\mathcal{L}/\partial\bm{b} = \bar{\bm{z}}.
  \label{eq:bp_afin}
\end{equation}
\end{proposicion}

\begin{proof}
$\partial\mathcal{L}/\partial a_k = \sum_j
(\partial\mathcal{L}/\partial z_j)(\partial z_j/\partial a_k) =
\sum_j \bar{z}_j W_{jk} = (\bm{W}^\top\bar{\bm{z}})_k$. Para los
pesos: $\partial\mathcal{L}/\partial W_{jk} = \bar{z}_j a_k =
(\bar{\bm{z}}\bm{a}^\top)_{jk}$. Para el sesgo:
$\partial z_j/\partial b_j = 1$. $\square$
\end{proof}

\begin{proposicion}{Backpropagation de la activación elemento a elemento}
\label{prop:bp_sigma}
Para $\bm{a} = \sigma(\bm{z})$ aplicada componente a componente:
\begin{equation}
  \bar{\bm{z}} = \bar{\bm{a}} \odot \sigma'(\bm{z}),
  \label{eq:bp_sigma}
\end{equation}
donde $\odot$ es el producto de Hadamard (elemento a elemento).
\end{proposicion}

\subsection{Las ecuaciones BP1–BP4}
\label{subsec:bp1234}

\begin{teorema}{Las cuatro ecuaciones de backpropagation}
\label{thm:cuatro_ec}
Para un MLP de $L$ capas con activaciones $\sigma^{(l)}$, definiendo el
\textbf{error de la capa} $\bm{\delta}^{(l)} =
\partial\mathcal{L}/\partial\bm{z}^{(l)}$:
\begin{align}
  \bm{\delta}^{(L)} &= \nabla_{\bm{a}^{(L)}}\ell
  \odot \sigma'^{(L)}(\bm{z}^{(L)})
  \tag{BP1}\label{eq:bp1}\\
  \bm{\delta}^{(l)} &= \bigl(\bm{W}^{(l+1)\top}
  \bm{\delta}^{(l+1)}\bigr) \odot \sigma'^{(l)}(\bm{z}^{(l)}),
  \quad l = L{-}1,\ldots,1
  \tag{BP2}\label{eq:bp2}\\
  \frac{\partial\mathcal{L}}{\partial\bm{W}^{(l)}}
  &= \bm{\delta}^{(l)}{\bm{a}^{(l-1)}}^\top
  \tag{BP3}\label{eq:bp3}\\
  \frac{\partial\mathcal{L}}{\partial\bm{b}^{(l)}}
  &= \bm{\delta}^{(l)}
  \tag{BP4}\label{eq:bp4}
\end{align}
\end{teorema}

\begin{proof}
\textbf{BP1:} por la regla de la cadena,
$\delta_j^{(L)} = \sum_k (\partial\ell/\partial a_k^{(L)})
(\partial a_k^{(L)}/\partial z_j^{(L)}) =
(\partial\ell/\partial a_j^{(L)})\sigma'^{(L)}(z_j^{(L)})$.
Para la capa de salida lineal (regresión):
$\sigma'^{(L)} = 1$ y $\bm{\delta}^{(L)} = \nabla_{\bm{a}^{(L)}}\ell$.

\textbf{BP2:} Por la regla de la cadena:
\begin{align*}
  \delta_j^{(l)}
  &= \frac{\partial\mathcal{L}}{\partial z_j^{(l)}}
  = \sum_k \frac{\partial\mathcal{L}}{\partial z_k^{(l+1)}}
  \frac{\partial z_k^{(l+1)}}{\partial a_j^{(l)}}
  \frac{\partial a_j^{(l)}}{\partial z_j^{(l)}} \\
  &= \sum_k \delta_k^{(l+1)}\, W_{kj}^{(l+1)}\,
  \sigma'^{(l)}(z_j^{(l)})
  = \bigl({\bm{W}^{(l+1)}}^\top\bm{\delta}^{(l+1)}\bigr)_j
  \cdot \sigma'^{(l)}(z_j^{(l)}).
\end{align*}

\textbf{BP3–BP4:} consecuencia de la
Proposición~\ref{prop:bp_afin} con $\bar{\bm{z}} = \bm{\delta}^{(l)}$,
$\bm{a} = \bm{a}^{(l-1)}$. $\square$
\end{proof}

\subsection{El algoritmo completo: forward y backward pass}
\label{subsec:alg_completo}

\begin{algorithm}[H]
\caption{Forward y backward pass del MLP}
\label{alg:backprop}
\KwIn{Entrada $\bm{x}$, etiqueta $y$, parámetros
  $\{\bm{W}^{(l)},\bm{b}^{(l)}\}_{l=1}^L$}
\KwOut{Pérdida $\mathcal{L}$, gradientes
  $\{\partial\mathcal{L}/\partial\bm{W}^{(l)},
  \partial\mathcal{L}/\partial\bm{b}^{(l)}\}_{l=1}^L$}

\tcp{--- Forward pass: almacenar activaciones ---}
$\bm{a}^{(0)} \leftarrow \bm{x}$\;
\For{$l \leftarrow 1$ \KwTo $L$}{
  $\bm{z}^{(l)} \leftarrow \bm{W}^{(l)}\bm{a}^{(l-1)}
  + \bm{b}^{(l)}$\;
  $\bm{a}^{(l)} \leftarrow \sigma^{(l)}(\bm{z}^{(l)})$\;
}
$\mathcal{L} \leftarrow \ell(\bm{a}^{(L)}, y)$\;

\tcp{--- Backward pass: propagar errores ---}
$\bm{\delta}^{(L)} \leftarrow \nabla_{\bm{a}^{(L)}}\ell
\odot \sigma'^{(L)}(\bm{z}^{(L)})$ \tcp*{BP1}
\For{$l \leftarrow L$ \KwTo $1$}{
  $\partial\mathcal{L}/\partial\bm{W}^{(l)} \leftarrow
  \bm{\delta}^{(l)}{\bm{a}^{(l-1)}}^\top$ \tcp*{BP3}
  $\partial\mathcal{L}/\partial\bm{b}^{(l)} \leftarrow
  \bm{\delta}^{(l)}$ \tcp*{BP4}
  \lIf{$l > 1$}{
    $\bm{\delta}^{(l-1)} \leftarrow
    (\bm{W}^{(l)\top}\bm{\delta}^{(l)}) \odot
    \sigma'^{(l-1)}(\bm{z}^{(l-1)})$ \tcp*{BP2}
  }
}
\Return $\mathcal{L}$,
  $\{\partial\mathcal{L}/\partial\bm{W}^{(l)},
  \partial\mathcal{L}/\partial\bm{b}^{(l)}\}_{l=1}^L$\;
\end{algorithm}

\begin{observacion}{Costo computacional}
El costo dominante del backward pass son los productos matriciales en BP2
($\bm{W}^{(l)\top}\bm{\delta}^{(l)}$) y BP3
($\bm{\delta}^{(l)}{\bm{a}^{(l-1)}}^\top$), de las mismas dimensiones que
los del forward pass. El costo total del backward es $\sim 2\times$ el
forward. Un ciclo completo (forward + backward) cuesta $\sim 3\times$ un
forward solo.
\end{observacion}

% ─────────────────────────────────────────────────────────────────────────────
\section{Backpropagation como método adjunto}
\label{sec:adjunto}
% ─────────────────────────────────────────────────────────────────────────────

\dificultad{3}{4}

Esta sección revela la estructura profunda del algoritmo: backpropagation
es el \textbf{método adjunto} aplicado al sistema dinámico discreto
definido por el forward pass del MLP.

\subsection{El MLP como sistema dinámico con restricciones}
\label{subsec:mlp_dinamico}

El forward pass define un sistema dinámico discreto:
\begin{equation}
  \bm{a}^{(l)} = \sigma(\bm{W}^{(l)}\bm{a}^{(l-1)}+\bm{b}^{(l)}),
  \qquad l = 1, \ldots, L,
  \label{eq:forward_dyn}
\end{equation}
con ``condición inicial'' $\bm{a}^{(0)} = \bm{x}$. El problema de
entrenamiento es minimizar
$\mathcal{L} = \ell(\bm{a}^{(L)}, y)$ sujeto a las
restricciones~\eqref{eq:forward_dyn}. Esto es un problema de
\textbf{control óptimo discreto}: las capas son instantes temporales, los
parámetros $(\bm{W}^{(l)}, \bm{b}^{(l)})$ son los controles, y las
activaciones $\bm{a}^{(l)}$ son el estado.

\subsection{Multiplicadores de Lagrange y la ecuación adjunta}
\label{subsec:lagrange}

Se introduce un multiplicador de Lagrange $\bm{p}^{(l)}$ para cada
restricción~\eqref{eq:forward_dyn} y se forma el lagrangiano:
\begin{equation}
  \mathfrak{L} = \ell(\bm{a}^{(L)},y) + \sum_{l=1}^L
  {\bm{p}^{(l)}}^\top\bigl[\sigma(\bm{W}^{(l)}\bm{a}^{(l-1)}
  +\bm{b}^{(l)}) - \bm{a}^{(l)}\bigr].
  \label{eq:lagrangiano}
\end{equation}
La condición de estacionariedad
$\partial\mathfrak{L}/\partial\bm{a}^{(l)} = 0$ para $l < L$ da:
\begin{equation}
  \bm{p}^{(l)} = \bm{W}^{(l+1)\top}\bigl(\bm{p}^{(l+1)} \odot
  \sigma'(\bm{z}^{(l+1)})\bigr),
  \label{eq:adjunta_discreta}
\end{equation}
con condición terminal $\bm{p}^{(L)} = \nabla_{\bm{a}^{(L)}}\ell$. Esto
es exactamente BP2 con la identificación
$\bm{\delta}^{(l)} = \bm{p}^{(l)} \odot \sigma'(\bm{z}^{(l)})$. El
backward pass es la \textbf{integración hacia atrás de la ecuación adjunta
discreta}, con condición terminal dada por BP1.

\subsection{El principio de Pontryagin como marco unificador}
\label{subsec:pontryagin}

En teoría de control óptimo, el principio del mínimo de Pontryagin
establece que la trayectoria óptima de un sistema dinámico con controles
satisface un sistema de dos ecuaciones: la ecuación del estado (forward) y
la ecuación adjunta (backward), acopladas por la condición de optimalidad
sobre los controles. El MLP es un caso particular con dinámica discreta y
controles no acotados:

\begin{center}
\begin{tabular}{ll}
\toprule
\textbf{Control óptimo} & \textbf{Backpropagation} \\
\midrule
Estado $\bm{x}(t)$ & Activación $\bm{a}^{(l)}$ \\
Control $\bm{u}(t)$ & Parámetros $(\bm{W}^{(l)}, \bm{b}^{(l)})$ \\
Dinámica $\dot{\bm{x}} = f(\bm{x},\bm{u})$ & Forward pass~\eqref{eq:forward_dyn} \\
Costo terminal $\ell(\bm{x}(T))$ & Pérdida $\ell(\bm{a}^{(L)}, y)$ \\
Co-estado $\bm{p}(t)$ & Error $\bm{\delta}^{(l)}$ \\
Ec.\ adjunta $\dot{\bm{p}} = -(\partial f/\partial\bm{x})^\top\bm{p}$ & BP2 \\
Cond.\ terminal $\bm{p}(T) = \nabla\ell$ & BP1 \\
Gradiente $\partial\ell/\partial\bm{u} = \bm{p}^\top\partial f/\partial\bm{u}$ & BP3 \\
\bottomrule
\end{tabular}
\end{center}

La profundidad $L$ no es solo un hiperparámetro arquitectónico: es el
\textbf{horizonte temporal} del sistema dinámico. El
Capítulo~\ref{ch:cap9} hará esta conexión explícita con las Neural ODEs,
donde la profundidad se vuelve continua.

\subsection{El método adjunto en la física computacional}
\label{subsec:adjunto_fisica}

El método adjunto no es un concepto de ML importado a la física. Es al
revés: la física lo inventó y el ML lo redescubrió. Cuatro ejemplos que
todo físico puede reconocer:

\textbf{Mecánica analítica.}
La formulación de Hamilton del oscilador armónico define el estado
$(\bm{q}, \bm{p})$ (posición generalizada y momento conjugado). Las
ecuaciones de Hamilton $\dot{\bm{q}} = \partial H/\partial\bm{p}$,
$\dot{\bm{p}} = -\partial H/\partial\bm{q}$ son formalmente idénticas al
par forward/backward: $\bm{q}$ es el estado que evoluciona hacia adelante
y $\bm{p}$ es el co-estado que propaga la información de gradiente. En
backpropagation, la activación $\bm{a}^{(l)}$ es la coordenada
generalizada y el error $\bm{\delta}^{(l)}$ es el momento conjugado.

\textbf{Óptica y el principio de Fermat.}
Un rayo de luz que viaja de un punto $A$ a un punto $B$ en un medio con
índice de refracción variable $n(\bm{r})$ sigue el camino que minimiza
el tiempo óptico $\mathcal{L} = \int_A^B n(\bm{r})\,ds$. La ecuación
del rayo $\frac{d}{ds}(n\hat{\bm{t}}) = \nabla n$ es la ecuación del
estado; la ecuación adjunta propaga la sensibilidad del tiempo óptico
respecto a perturbaciones del índice. En diseño óptico, calcular cómo
cambia la aberración de una lente ante perturbaciones de curvatura es un
problema adjunto: el backward pass a través del sistema óptico.

\textbf{Problemas inversos en espectroscopía.}
En la deconvolución de un espectro medido
$y(\omega) = \int h(\omega-\omega')s(\omega')\,d\omega' + n(\omega)$, el
operador forward es la convolución con la función de respuesta $h$. El
operador adjunto es la convolución con $h(-\omega)$ (correlación). La
inversión por el método adjunto (iteración de Landweber) es:
$s^{(k+1)} = s^{(k)} + \eta\, h * (y - h * s^{(k)})$, donde $h *$
denota la convolución adjunta. Esto es formalmente GD con el gradiente
calculado por el operador adjunto: exactamente backpropagation sobre un
``MLP de una capa'' cuyo peso es $h$.

\textbf{Sismología inversa (Tarantola, 1984).}
La inversión de forma de onda completa (\textit{full waveform inversion})
calcula el gradiente de la función de misfit respecto a las velocidades
sísmicas del subsuelo propagando el campo de onda \emph{hacia atrás} en
el tiempo (backward wavefield). Tarantola
(1984)~\cite{tarantola1984inversion} identificó esta propagación inversa
como el método adjunto aplicado a la ecuación de onda, antes de que la
comunidad de ML reconociera la misma estructura en backpropagation. La
conexión fue redescubierta formalmente por Griewank
(1989)~\cite{griewank2008evaluating} en el contexto de la diferenciación
automática.

% ─────────────────────────────────────────────────────────────────────────────
\section{Extensión al caso de mini-lote}
\label{sec:bp_batch}
% ─────────────────────────────────────────────────────────────────────────────

\dificultad{2}{3}

Para un mini-lote $\mathcal{B}$ de tamaño $B$, el forward pass produce
matrices $\bm{A}^{(l)} \in \mathbb{R}^{B \times n_l}$. Las ecuaciones BP
se extienden naturalmente:
\begin{align}
  \bm{\Delta}^{(L)} &= \nabla_{\bm{A}^{(L)}}\mathcal{L}_{\mathcal{B}}
  \odot \sigma'^{(L)}(\bm{Z}^{(L)}) \tag{BP1-B}\\
  \bm{\Delta}^{(l)} &= (\bm{\Delta}^{(l+1)}\bm{W}^{(l+1)})
  \odot \sigma'^{(l)}(\bm{Z}^{(l)}) \tag{BP2-B}\\
  \frac{\partial\mathcal{L}_{\mathcal{B}}}{\partial\bm{W}^{(l)}}
  &= \frac{1}{B}\bm{\Delta}^{(l)\top}\bm{A}^{(l-1)} \tag{BP3-B}\\
  \frac{\partial\mathcal{L}_{\mathcal{B}}}{\partial\bm{b}^{(l)}}
  &= \frac{1}{B}\bm{\Delta}^{(l)\top}\bm{1}_B \tag{BP4-B}
\end{align}
El costo dominante es $O(Bn_ln_{l-1})$ por capa: el mismo que el forward.

% ─────────────────────────────────────────────────────────────────────────────
\section{Reglas de backpropagation para operaciones estándar}
\label{sec:tabla_bp}
% ─────────────────────────────────────────────────────────────────────────────

\dificultad{2}{3}

\subsection{Softmax y entropía cruzada}
\label{subsec:bp_softmax}

Para $\hat{p}_j = e^{z_j}/\sum_k e^{z_k}$ (softmax) seguida de
$\ell = -\sum_j y_j\log\hat{p}_j$ (entropía cruzada), el gradiente
combinado se simplifica a:
\begin{equation}
  \partial\mathcal{L}/\partial\bm{z}^{(L)} = \hat{\bm{p}} - \bm{y}.
  \label{eq:bp_softmax_ce}
\end{equation}
La derivación directa evita formar la jacobiana del softmax ($O(C^2)$ con
$C$ clases) y es numéricamente estable.

\subsection{Batch normalization}
\label{subsec:bp_bn}

Para la normalización por lotes $\hat{z}_j^{(i)} = (z_j^{(i)} -
\mu_j)/\sqrt{s_j^2+\varepsilon}$ seguida de la transformación afín
$\tilde{z}_j^{(i)} = \gamma_j\hat{z}_j^{(i)} + \beta_j$, el gradiente
respecto a $\bm{z}$ involucra las estadísticas del batch ($\mu_j$,
$s_j^2$) y no es diagonal:
\begin{equation}
  \frac{\partial\mathcal{L}}{\partial z_j^{(i)}}
  = \frac{\gamma_j}{B\sqrt{s_j^2+\varepsilon}}
  \left[B\bar{\hat{z}}_j^{(i)} - \sum_k\bar{\hat{z}}_j^{(k)}
  - \hat{z}_j^{(i)}\sum_k\bar{\hat{z}}_j^{(k)}\hat{z}_j^{(k)}\right].
  \label{eq:bp_bn}
\end{equation}
El Ejercicio Resuelto~\ref{ej:bp_batchnorm} deriva esta expresión.

\subsection{Conexiones de salto y concatenación}
\label{subsec:bp_skip}

Para la suma de dos ramas $\bm{y} = \bm{x}_1 + \bm{x}_2$:
$\bar{\bm{x}}_1 = \bar{\bm{y}}$, $\bar{\bm{x}}_2 = \bar{\bm{y}}$.
El gradiente se distribuye sin modificación a ambas ramas. Esta es la
razón por la que las conexiones de salto (ResNet) mitigan el problema del
gradiente que se desvanece: el gradiente tiene un camino directo que no
pasa por las activaciones $\sigma'$.

Para la concatenación $\bm{y} = [\bm{x}_1;\bm{x}_2]$: el gradiente se
divide según las componentes.

% ─────────────────────────────────────────────────────────────────────────────
\section{Implementación y gradient checkpointing}
\label{sec:implementacion}
% ─────────────────────────────────────────────────────────────────────────────

\dificultad{2}{2}

\subsection{Backpropagation en frameworks modernos}
\label{subsec:frameworks}

En la práctica, las ecuaciones BP1–BP4 no se implementan manualmente. Los
frameworks (PyTorch, JAX) construyen el grafo computacional durante el
forward pass y ejecutan el backward automáticamente. Entender las
ecuaciones es imprescindible para diagnosticar problemas (gradientes nulos
o explosivos), implementar operaciones personalizadas, y comprender qué
calcula el framework.

\begin{verbatim}
# PyTorch: backward automático
import torch, torch.nn as nn

modelo = nn.Sequential(
    nn.Linear(1, 32), nn.ReLU(),
    nn.Linear(32, 1)
)
x = torch.randn(25, 1)
y = torch.sin(2 * 3.14159 * x) + 0.15 * torch.randn(25, 1)

pred = modelo(x)                       # Forward
loss = nn.MSELoss()(pred, y)           # Pérdida
loss.backward()                        # BP1-BP4 automático

for nombre, param in modelo.named_parameters():
    print(nombre, param.grad.shape)    # dL/dW, dL/db
\end{verbatim}

\begin{verbatim}
# JAX: diferenciación funcional
import jax, jax.numpy as jnp

def perdida(params, x, y):
    a = x
    for W, b in params[:-1]:
        a = jax.nn.relu(a @ W.T + b)
    W, b = params[-1]
    return jnp.mean((a @ W.T + b - y) ** 2)

grad_fn = jax.grad(perdida)            # Transforma la función
grads = grad_fn(params, x, y)          # Aplica BP1-BP4
\end{verbatim}

En ambos ejemplos, la red es el MLP $(1, 32, 1)$ del oscilador del
Capítulo~\ref{ch:cap4}: los datos son $\sin(2\pi t)$ con ruido.

\subsection{Gradient checkpointing}
\label{subsec:checkpointing}

El backward pass requiere las activaciones
$\{\bm{z}^{(l)}, \bm{a}^{(l)}\}$ del forward. Para un modelo con
$L = 100$ capas, anchura $m = 1024$ y batch $B = 512$ en 32 bits:
$\text{memoria} = 2LmB \times 4\,\text{bytes} \approx 419\,\text{MB}$.
Para modelos de lenguaje ($L = 96$, $m = 12288$, $B = 2048$):
$\sim 100\,\text{GB}$.

El \textbf{gradient checkpointing} almacena solo $\sqrt{L}$ activaciones
en puntos igualmente espaciados y recomputa las intermedias durante el
backward. La memoria se reduce de $O(LmB)$ a $O(\sqrt{L}\,mB)$ a costa
de $\sim 2\times$ cómputo adicional. Es el intercambio estándar
memoria/cómputo en modelos grandes.

% ─────────────────────────────────────────────────────────────────────────────
\section{El oscilador como laboratorio: backpropagation manual}
\label{sec:oscilador6}
% ─────────────────────────────────────────────────────────────────────────────

\dificultad{2}{2}

Retomamos el MLP $(1, 32, 1)$ con ReLU entrenado sobre los datos del
oscilador. Para visualizar qué calcula backpropagation, reducimos la
arquitectura a $(1, 3, 1)$ (tres neuronas ocultas) y ejecutamos el
algoritmo con valores numéricos.

\textbf{Parámetros:}
$\bm{W}^{(1)} = (2, -1, 0.5)^\top$,
$\bm{b}^{(1)} = (-1, 0.5, -0.3)^\top$,
$\bm{W}^{(2)} = (1, -1, 2)$, $b^{(2)} = 0$.
Entrada: $x = 0.25$ (un instante del oscilador).
Dato: $y = \sin(2\pi \times 0.25) = 1$.

\textbf{Forward pass:}
$\bm{z}^{(1)} = \bm{W}^{(1)}x + \bm{b}^{(1)}
= (2(0.25)-1,\; -0.25+0.5,\; 0.125-0.3)^\top = (-0.5, 0.25, -0.175)^\top$.
$\bm{a}^{(1)} = \text{ReLU}(\bm{z}^{(1)}) = (0, 0.25, 0)^\top$.
$z^{(2)} = \bm{W}^{(2)}\bm{a}^{(1)} + b^{(2)} = 0 - 0.25 + 0 = -0.25$.
$\hat{y} = -0.25$. Pérdida: $\mathcal{L} = (-0.25 - 1)^2 = 1.5625$.

\textbf{Backward pass (BP1–BP4):}

BP1: $\delta^{(2)} = 2(\hat{y}-y) = 2(-0.25-1) = -2.5$ (para MSE con
capa de salida lineal).

BP2: $\bm{\delta}^{(1)} = (\bm{W}^{(2)\top}\delta^{(2)}) \odot
\text{ReLU}'(\bm{z}^{(1)}) = (-2.5, 2.5, -5)^\top \odot (0, 1, 0)^\top
= (0, 2.5, 0)^\top$.
Las neuronas 1 y 3 tenían preactivación negativa (inactivas): su
gradiente es cero. Solo la neurona 2 propaga gradiente.

BP3: $\partial\mathcal{L}/\partial\bm{W}^{(2)} =
\delta^{(2)}{\bm{a}^{(1)}}^\top = -2.5 \times (0, 0.25, 0)
= (0, -0.625, 0)$.
$\partial\mathcal{L}/\partial\bm{W}^{(1)} =
\bm{\delta}^{(1)} x = (0, 2.5, 0)^\top \times 0.25
= (0, 0.625, 0)^\top$.

BP4: $\partial\mathcal{L}/\partial b^{(2)} = -2.5$.
$\partial\mathcal{L}/\partial\bm{b}^{(1)} = (0, 2.5, 0)^\top$.

Solo los parámetros de la neurona activa (neurona 2) reciben
actualización. Las neuronas inactivas tienen gradiente cero: el SGD no
las mueve. Esta observación tiene una consecuencia para el entrenamiento:
si una neurona ReLU se inactiva para \emph{todas} las muestras del
dataset (``neurona muerta''), nunca se recupera. La inicialización de He
(Capítulo~\ref{ch:cap5}) minimiza este riesgo.

El Ejercicio Computacional~\ref{ejcomp:oscilador6} verifica estos cálculos
contra diferencias finitas y extiende el análisis a la red $(1, 32, 1)$
completa.

% ─────────────────────────────────────────────────────────────────────────────
\section{Resumen del Capítulo 6 y transición}
\label{sec:resumen6}
% ─────────────────────────────────────────────────────────────────────────────

Este capítulo respondió la pregunta que el Capítulo~\ref{ch:cap5} dejó
abierta: cómo calcular $\nabla_{\bm{\theta}}\mathcal{L}$ eficientemente.

A nivel algorítmico, backpropagation es el modo reverso de la
diferenciación automática sobre el grafo computacional. Las ecuaciones
BP1–BP4 propagan el gradiente desde la salida hacia la entrada,
acumulando los gradientes de los parámetros en cada capa. El costo es
$\sim 3\times$ un forward pass, independiente de $p$.

A nivel estructural, backpropagation es el método adjunto aplicado al
sistema dinámico discreto definido por el MLP. Los errores
$\bm{\delta}^{(l)}$ son los multiplicadores de Lagrange del problema de
control óptimo con las capas como instantes temporales. El principio de
Pontryagin es la versión continua; las Neural ODEs del
Capítulo~\ref{ch:cap9} harán esta conexión explícita.

Los Capítulos~\ref{ch:cap4}, \ref{ch:cap5} y \ref{ch:cap6} cierran el
ciclo fundamental: el objeto (MLP), el optimizador (SGD/Adam) y el
gradiente (backpropagation). El Capítulo~\ref{ch:cap7} cambia de
perspectiva: en lugar de preguntar ``¿cuál es el mejor $\bm{\theta}^*$?''
(estimación puntual), pregunta ``¿cuál es la distribución
$p(\bm{\theta}|\mathcal{D})$?'' (inferencia bayesiana). La respuesta
revelará que la pérdida MSE es el logaritmo de una verosimilitud
gaussiana, la regularización $\ell_2$ del Capítulo~\ref{ch:cap2} es un
prior gaussiano, y la distribución de Gibbs del
Capítulo~\ref{ch:cap5} es la distribución de máxima entropía. Cada pieza
que se construyó pragmáticamente tiene una justificación probabilística
precisa.

% ─────────────────────────────────────────────────────────────────────────────
\section*{Ejercicios resueltos}
\addcontentsline{toc}{section}{Ejercicios resueltos}
% ─────────────────────────────────────────────────────────────────────────────

\begin{ejercicio}
\label{ej:bp_manual}
\textbf{(Backpropagation manual en el MLP del oscilador).}

Verificar los cálculos de la Sección~\ref{sec:oscilador6} para el MLP
$(1, 3, 1)$ con los parámetros dados, entrada $x = 0.5$ y dato
$y = \sin(2\pi \times 0.5) = 0$.

\medskip\noindent\textit{Solución.}
$\bm{z}^{(1)} = (2(0.5)-1, -0.5+0.5, 0.25-0.3)^\top = (0, 0, -0.05)^\top$.
$\bm{a}^{(1)} = \text{ReLU}(\bm{z}^{(1)}) = (0, 0, 0)^\top$.
$z^{(2)} = 0$. $\hat{y} = 0$. $\mathcal{L} = (0-0)^2 = 0$.

BP1: $\delta^{(2)} = 2(0-0) = 0$.

Todos los gradientes son cero: la pérdida es exactamente cero para
este dato (por coincidencia: $\sin(\pi) = 0$ y la red predice $0$).
No se produce ninguna actualización. $\square$
\end{ejercicio}

\begin{ejercicio}
\label{ej:bp_batchnorm}
\textbf{(Backpropagation a través de batch normalization).}

Derivar la ecuación~\eqref{eq:bp_bn} para el gradiente de la
pérdida respecto a las preactivaciones $z_j^{(i)}$ a través de la
capa de batch normalization.

\medskip\noindent\textit{Solución.}
Sea $\hat{z}_j^{(i)} = (z_j^{(i)}-\mu_j)/\sigma_j$ con
$\mu_j = \frac{1}{B}\sum_k z_j^{(k)}$ y
$\sigma_j = \sqrt{s_j^2+\varepsilon}$.

El gradiente $\partial\mathcal{L}/\partial z_j^{(i)}$ tiene tres
contribuciones (por la dependencia de $\hat{z}_j^{(i)}$, $\mu_j$ y
$s_j^2$ en $z_j^{(i)}$):
\begin{align*}
  \frac{\partial\mathcal{L}}{\partial z_j^{(i)}}
  &= \frac{\partial\mathcal{L}}{\partial\hat{z}_j^{(i)}}
  \frac{1}{\sigma_j}
  + \frac{\partial\mathcal{L}}{\partial\mu_j}\frac{1}{B}
  + \frac{\partial\mathcal{L}}{\partial s_j^2}
  \frac{2(z_j^{(i)}-\mu_j)}{B}.
\end{align*}
Con $\bar{\hat{z}}_j^{(i)} = \gamma_j
\partial\mathcal{L}/\partial\tilde{z}_j^{(i)}$ y calculando las
derivadas respecto a $\mu_j$ y $s_j^2$, se llega
a~\eqref{eq:bp_bn}. $\square$
\end{ejercicio}

% ─────────────────────────────────────────────────────────────────────────────
\section*{Ejercicios propuestos}
\addcontentsline{toc}{section}{Ejercicios propuestos}
% ─────────────────────────────────────────────────────────────────────────────

\begin{ejercicio}
\textbf{(Verificación del gradiente por diferencias finitas).}
\begin{enumerate}[label=(\alph*)]
  \item Implementar forward y backward manual del
        Algoritmo~\ref{alg:backprop} para el MLP $(1, 3, 1)$ con ReLU.
  \item Para parámetros y entrada del oscilador ($x = 0.25$, $y = 1$),
        calcular $\partial\mathcal{L}/\partial\theta_j$ por diferencias
        finitas centradas ($h = 10^{-5}$) para 5 parámetros.
  \item Comparar con BP. Error relativo esperado: $O(h^2) \sim 10^{-10}$.
  \item ¿Qué ocurre con $h = 10^{-1}$ (truncación) y $h = 10^{-12}$
        (cancelación catastrófica)?
\end{enumerate}
\end{ejercicio}

\begin{ejercicio}
\textbf{(Softmax + entropía cruzada).}
\begin{enumerate}[label=(\alph*)]
  \item Calcular el jacobiano del softmax:
        $\partial\hat{p}_j/\partial z_k = \hat{p}_j(\delta_{jk}-\hat{p}_k)$.
  \item Combinar con $\partial\ell/\partial\hat{p}_j = -y_j/\hat{p}_j$
        para obtener~\eqref{eq:bp_softmax_ce}.
  \item ¿Por qué los frameworks implementan la versión combinada en lugar
        de aplicar el jacobiano del softmax por separado?
\end{enumerate}
\end{ejercicio}

\begin{ejercicio}
\textbf{(Conexiones de salto y el gradiente).}
Red residual: $\bm{a}^{(2)} = \bm{a}^{(0)} + F(\bm{a}^{(0)};\bm{W}^{(1)})$,
$\bm{a}^{(4)} = \bm{a}^{(2)} + F(\bm{a}^{(2)};\bm{W}^{(2)})$, con
$F(\bm{a};\bm{W}) = \text{ReLU}(\bm{W}\bm{a})$.
\begin{enumerate}[label=(\alph*)]
  \item Dibujar el grafo computacional.
  \item Calcular $\partial\mathcal{L}/\partial\bm{a}^{(0)}$ por el modo
        reverso. ¿Cuántos caminos conectan $\mathcal{L}$ con $\bm{a}^{(0)}$?
  \item Mostrar que el gradiente incluye un término identidad.
        ¿Qué implica para el gradiente que se desvanece?
\end{enumerate}
\end{ejercicio}

% ─────────────────────────────────────────────────────────────────────────────
\section*{Ejercicios computacionales}
\addcontentsline{toc}{section}{Ejercicios computacionales}
% ─────────────────────────────────────────────────────────────────────────────

\begin{ejerciciocomp}{El oscilador: backpropagation desde cero vs.\ PyTorch}
\label{ejcomp:oscilador6}
\textbf{Objetivo:} implementar el Algoritmo~\ref{alg:backprop} desde cero
en NumPy para el MLP del oscilador y verificar contra PyTorch/JAX y contra
diferencias finitas.

\medskip\noindent\textbf{Datos:}
MLP $(1, 32, 1)$ con ReLU, inicialización de He, $N = 25$ puntos del
oscilador $\sin(2\pi t) + \varepsilon_i$, $\sigma = 0.15$.

\medskip\noindent\textbf{Algoritmo:}
\begin{enumerate}
  \item Implementar \texttt{forward(params, x)} que devuelva la salida y
        almacene $\{(\bm{z}^{(l)}, \bm{a}^{(l)})\}$.
  \item Implementar \texttt{backward(cache, y)} que aplique BP1–BP4 y
        devuelva los gradientes.
  \item Para una muestra $(t_1, y_1) = (0.25, \sin(\pi/2)+\varepsilon_1)$:
        comparar los gradientes de la implementación manual con
        \texttt{loss.backward()} de PyTorch. Error relativo esperado:
        $< 10^{-6}$.
  \item Comparar con diferencias finitas centradas ($h = 10^{-5}$) para
        10 parámetros elegidos al azar. Error esperado: $< 10^{-7}$.
  \item Ejecutar 500 iteraciones de SGD ($\eta = 10^{-3}$, $B = 4$) usando
        la implementación manual. Verificar que
        $\mathcal{L}_{\text{train}}$ decrece.
\end{enumerate}

\medskip\noindent\textbf{Verificación:}
Error relativo BP vs.\ PyTorch: $< 10^{-6}$ (diferencias solo por orden
de operaciones flotantes).
Error relativo BP vs.\ diferencias finitas: $< 10^{-7}$ (errores mayores
indican un bug).
Después de 500 iteraciones: $\mathcal{L}_{\text{train}} < 0.1$.
\end{ejerciciocomp}

\begin{ejerciciocomp}{Flujo de gradientes en profundidad}
\label{ejcomp:vanishing}
\textbf{Objetivo:} medir empíricamente el gradiente que se desvanece y
verificar que la inicialización de He lo resuelve para ReLU.

\medskip\noindent\textbf{Algoritmo:}
\begin{enumerate}
  \item MLPs de profundidad $L = 5, 10, 20$, anchura $m = 64$, activación
        sigmoide, inicialización $\mathcal{N}(0, 1/m)$.
  \item Para entrada y pérdida aleatorias, registrar
        $\|\partial\mathcal{L}/\partial\bm{W}^{(l)}\|_F$ por capa.
  \item Graficar
        $\log\|\partial\mathcal{L}/\partial\bm{W}^{(l)}\|_F$ vs.\ $l$.
        ¿La pendiente es consistente con atenuación exponencial?
  \item Repetir con ReLU + He. ¿Desaparece la atenuación?
\end{enumerate}

\medskip\noindent\textbf{Verificación:}
Sigmoide: pendiente $\approx -0.15$ por capa (factor $\sim 0.86$ por capa,
consistente con $\max|\sigma'| = 0.25$ y $m\sigma_w^2 = 1$).
ReLU + He: pendiente $\approx 0$ (gradiente estable).
\end{ejerciciocomp}

% ─────────────────────────────────────────────────────────────────────────────
\begin{notasbib}
El algoritmo de backpropagation fue redescubierto independientemente varias
veces. Las referencias más influyentes son Werbos
(1974)~\cite{werbos1974beyond}, que lo derivó en el contexto del control
óptimo en su tesis doctoral, y Rumelhart, Hinton y Williams
(1986)~\cite{rumelhart1986learning}, que lo popularizaron demostrando su
eficacia para aprender representaciones internas. La conexión con el
principio de Pontryagin~\cite{pontryagin1962mathematical} es bien conocida
en la comunidad de control; la formulación moderna del método adjunto
continuo puede encontrarse en Bryson y Ho
(1975)~\cite{bryson1975applied}.

La diferenciación automática como disciplina fue formalizada por Griewank
y Walther (2008)~\cite{griewank2008evaluating}. Los dos modos y sus
complejidades relativas están tratados en detalle en ese texto.

La conexión entre backpropagation y el método adjunto en sismología fue
establecida por Tarantola (1984)~\cite{tarantola1984inversion} en la
inversión de forma de onda completa. El uso del operador adjunto en
problemas inversos de espectroscopía y óptica es estándar; véase
Hansen (1998)~\cite{hansen1998rank} para una perspectiva de problemas
inversos y Vogel (2002)~\cite{vogel2002computational} para métodos
adjuntos en problemas inversos computacionales.

El gradient checkpointing fue propuesto por Chen et al.\
(2016)~\cite{chen2016training} con la cota $O(\sqrt{L})$. La rediscusión
en el contexto de Neural ODEs por Chen et al.\
(2018)~\cite{chen2018neural} cerró el círculo: el método adjunto continuo
de Pontryagin aplicado al MLP como ODE discreta.

La batch normalization: Ioffe y Szegedy
(2015)~\cite{ioffe2015batch}.
\end{notasbib}

% ─────────────────────────────────────────────────────────────────────────────
% ESTRELLA POLAR — Sección independiente
% ─────────────────────────────────────────────────────────────────────────────

\begin{estrella}{El método adjunto en el sustituto neutrónico}
\label{sec:estrella6}

\textbf{El problema de la cinética inversa diferenciable.}
El sustituto MLP del reactor (Capítulos~\ref{ch:cap4}--\ref{ch:cap5})
aproxima $f_{\text{MLP}}(\bm{u}) \approx \bm{\Phi}(\bm{u})$. En
análisis de seguridad se necesita el jacobiano
$\partial\bm{\Phi}/\partial\bm{u} \in \mathbb{R}^{K \times n}$ (cómo
cambia el flujo ante perturbaciones de los parámetros operacionales). Con
un código de Monte Carlo, calcular este jacobiano requeriría $n$
simulaciones adicionales (una por parámetro perturbado).

\textbf{El backward pass como solución.}
Con el sustituto diferenciable, el jacobiano se obtiene directamente.
Para $K = 500$ nodos de flujo y $n = 5$ parámetros operacionales, el modo
forward de AD es la elección correcta: $n = 5$ pasadas forward dan las 5
columnas del jacobiano, vs.\ $K = 500$ pasadas reverso para las filas. Es
la situación donde $n \ll K$ y el modo forward gana.

\textbf{Análisis de sensibilidad de seguridad.}
Las sensibilidades $\partial k_{\text{eff}}/\partial\sigma_x$ (cómo cambia
la reactividad ante perturbaciones en secciones eficaces) se obtienen
analíticamente via el método adjunto sobre el sustituto. El código TSUNAMI
(SCALE) implementa exactamente el mismo método sobre el operador de
transporte de Boltzmann: es el backward pass sobre la ecuación de difusión
neutrónica en lugar del grafo computacional del MLP.

\textbf{Conexión con el capítulo.}
BP2 es la ecuación adjunta discreta del forward pass. En el sustituto
neutrónico, propaga las sensibilidades del flujo desde los nodos de salida
hasta los parámetros de entrada: cada capa aplica $\bm{W}^{(l)\top}$ (el
operador adjunto de la transformación forward) ponderada por $\sigma'$.
Es textualmente el principio de Pontryagin sobre el sistema dinámico
discreto del MLP.
\end{estrella}

\printbibliography[heading=subbibliography]